\documentclass[answers,12pt,addpoints]{exam} 
\usepackage{import}

\import{C:/Users/prana/OneDrive/Desktop/MathNotes}{style.tex}

% Header
\newcommand{\name}{Pranav Tikkawar}
\newcommand{\course}{Stochastic Processes - Stats 654}
\newcommand{\assignment}{Reference Sheet}
\author{\name}
\title{\course \ - \assignment}

\begin{document}
\maketitle
\tableofcontents

\newpage

\section*{Problem-Solving Legend (What Tool Should I Try First?)}

\begin{itemize}
  \item \textbf{Step 0: Identify the object.}
  \begin{itemize}
    \item Is it about sums / random walks / hitting times? $\Rightarrow$ Think \emph{martingales + optional stopping}.
    \item Is it about long-run behavior / stationary distribution / convergence? $\Rightarrow$ Think \emph{Markov chains + mixing}.
    \item Is it about bounding $\mathbb{P}(\max \ge a)$ or tails? $\Rightarrow$ Think \emph{Doob / Kolmogorov / Azuma / McDiarmid}.
  \end{itemize}

  \item \textbf{Martingale tools (RW, hitting probs, expectations).}
  \begin{itemize}
    \item Try to build a martingale: $S_n$, $S_n-n\mu$, $S_n^2-n\sigma^2$, $\rho^{S_n}$ or a Doob martingale.
    \item If there is a natural random time (first hit, ruin, pattern), check if it is a stopping time and apply OST.
  \end{itemize}

  \item \textbf{Markov chain tools (stationarity, convergence).}
  \begin{itemize}
    \item Compute $\pi$ via $\pi P=\pi$ or symmetry/degree rule; check detailed balance for reversibility.
    \item Check irreducible + aperiodic; then use the convergence theorem or spectral gap/mixing bounds.
  \end{itemize}

  \item \textbf{Coupling tools (mixing bounds, distances).}
  \begin{itemize}
    \item Try a shared-uniform coupling so both chains use the same randomness.
    \item Bound $\mathbb{P}(\tau>t)$ (meeting time) and then use $\|P^t(x,\cdot)-\pi\|_{\mathrm{TV}}\le \mathbb{P}(\tau>t)$.
  \end{itemize}

  \item \textbf{Concentration tools (tail bounds).}
  \begin{itemize}
    \item Independent bounded variables $\Rightarrow$ Hoeffding / Chernoff.
    \item Martingale with bounded increments $\Rightarrow$ Azuma.
    \item Function of independent variables with bounded influence $\Rightarrow$ Doob martingale + McDiarmid.
  \end{itemize}

  \item \textbf{MCMC tools (MH and Gibbs).}
  \begin{itemize}
    \item MH: write proposal $Q$, compute $\alpha(x,y)=\min\big(1,\frac{\pi(y)Q(y,x)}{\pi(x)Q(x,y)}\big)$, then check detailed balance.
    \item Gibbs: derive full conditionals using the proportional trick, then argue invariance of $\pi$.
  \end{itemize}

  \item \textbf{Graph / RW tools (recurrence, structure).}
  \begin{itemize}
    \item Random walk on graph: treat as Markov chain; stationary $\pi(x)\propto\deg(x)$ for undirected graphs.
    \item On trees or higher-dimensional lattices, reduce to a (possibly biased) 1D distance process and reuse gambler’s ruin / biased RW results.
  \end{itemize}
\end{itemize}

\section{L1–L2: Conditional Expectation and Basic Martingales}

\begin{definition}[Conditional Expectation]
Let $(\Omega,\mathcal{F},\mathbb{P})$ be a probability space, $X$ an integrable random variable with $\mathbb{E}|X|<\infty$, and $\mathcal{G}\subseteq\mathcal{F}$ a sub-$\sigma$-algebra. A random variable $Y$ is called the \emph{conditional expectation} of $X$ given $\mathcal{G}$ if:
\begin{enumerate}
\item $Y$ is $\mathcal{G}$-measurable,
\item for every $A\in\mathcal{G}$,
\[
\int_A Y\,d\mathbb{P} = \int_A X\,d\mathbb{P}.
\]
\end{enumerate}
We write $Y = \mathbb{E}[X\mid \mathcal{G}]$. For a random variable $Z$, $\mathbb{E}[X\mid Z] := \mathbb{E}[X\mid \sigma(Z)]$.
\end{definition}

\begin{definition}[Filtration]
A \emph{filtration} on $(\Omega,\mathcal{F},\mathbb{P})$ is an increasing sequence of sub-$\sigma$-algebras
\[
\mathcal{F}_0 \subseteq \mathcal{F}_1 \subseteq \mathcal{F}_2 \subseteq \cdots \subseteq \mathcal{F}.
\]
We interpret $\mathcal{F}_n$ as the information available up to time $n$.
\end{definition}

\begin{definition}[Martingale, Submartingale, Supermartingale]
Let $(\mathcal{F}_n)_{n\ge 0}$ be a filtration. A stochastic process $(M_n)_{n\ge 0}$ is a \emph{martingale} w.r.t.\ $(\mathcal{F}_n)$ if:
\begin{enumerate}
\item $M_n$ is $\mathcal{F}_n$-measurable (adapted),
\item $\mathbb{E}|M_n|<\infty$ for every $n$,
\item $\mathbb{E}[M_{n+1}\mid \mathcal{F}_n] = M_n$ a.s.\ for all $n$.
\end{enumerate}
It is a \emph{submartingale} if $\mathbb{E}[M_{n+1}\mid \mathcal{F}_n]\ge M_n$ a.s., and a \emph{supermartingale} if $\mathbb{E}[M_{n+1}\mid \mathcal{F}_n]\le M_n$ a.s.
\end{definition}

\begin{example}[Classic Martingales from Random Walks]
Let $(X_i)_{i\ge 1}$ be i.i.d.\ with $\mathbb{E}X_i = \mu$ and $\mathrm{Var}(X_i)=\sigma^2<\infty$, and define $S_n = \sum_{i=1}^n X_i$ with natural filtration $\mathcal{F}_n = \sigma(X_1,\dots,X_n)$.
\begin{itemize}
\item If $\mu=0$, then $(S_n)$ is a martingale.
\item $(S_n - n\mu)$ is always a martingale.
\item $(S_n^2 - n\sigma^2)$ is a martingale.
\item For simple symmetric random walk ($X_i=\pm 1$), $S_n$ is a martingale and $S_n^2-n$ is a martingale.
\end{itemize}
\end{example}

\begin{definition}[Stopping Time]
Given a filtration $(\mathcal{F}_n)_{n\ge 0}$, a random variable $T:\Omega\to\{0,1,2,\dots\}\cup\{\infty\}$ is a \emph{stopping time} if
\[
\{T = n\}\in \mathcal{F}_n \quad \text{for all }n
\]
equivalently, $\{T\le n\}\in\mathcal{F}_n$ for all $n$. Intuitively, at time $n$ we can decide whether $T$ has occurred using only information up to $n$.
\end{definition}

\begin{definition}[Stopped Process]
Given a process $(M_n)$ and stopping time $T$, the \emph{stopped process} is
\[
M_{n\wedge T} := M_{\min(n,T)}.
\]
If $(M_n)$ is a martingale, then $(M_{n\wedge T})$ is also a martingale.
\end{definition}

\section{L5–L7: Optional Stopping, Convergence, Transformations}

\begin{theorem}[Optional Stopping Theorem (Bounded Case)]
Let $(M_n)_{n\ge 0}$ be a martingale w.r.t.\ $(\mathcal{F}_n)$, and let $T$ be a stopping time such that $T\le N$ a.s.\ for some deterministic $N$. Then
\[
\mathbb{E}[M_T] = \mathbb{E}[M_0].
\]
More generally, OST holds under suitable integrability and bounded-increment conditions.
\end{theorem}

\begin{theorem}[Wald's Identity (Version I)]
Let $(X_i)_{i\ge 1}$ be i.i.d.\ with $\mathbb{E}|X_1|<\infty$ and mean $\mu$, and let $T$ be an integrable stopping time such that $\mathbb{E}T<\infty$ and certain regularity conditions hold (e.g.\ $T$ bounded or suitable integrability). With $S_T = \sum_{i=1}^T X_i$, we have
\[
\mathbb{E}[S_T] = \mu\, \mathbb{E}[T].
\]
This follows by applying OST to the martingale $S_n - n\mu$.
\end{theorem}

\begin{definition}[Almost Sure and $L^1$ Convergence]
We say $X_n\to X$ almost surely (a.s.) if
\[
\mathbb{P}\big(\{\omega: X_n(\omega)\to X(\omega)\}\big)=1.
\]
We say $X_n\to X$ in expectation (or in $L^1$) if $\mathbb{E}|X_n-X|\to 0$.
\end{definition}

\begin{theorem}[Monotone Convergence Theorem]
If $(X_n)$ is a nondecreasing sequence of nonnegative random variables ($0\le X_1\le X_2\le\cdots$) with $X_n\to X$ a.s., then
\[
\lim_{n\to\infty}\mathbb{E}[X_n]=\mathbb{E}[X].
\]
\end{theorem}

\begin{theorem}[Dominated Convergence Theorem]
If $X_n\to X$ a.s.\ and there exists integrable $Y$ with $|X_n|\le Y$ a.s.\ for all $n$, then
\[
\lim_{n\to\infty}\mathbb{E}[X_n]=\mathbb{E}[X].
\]
\end{theorem}

\begin{lemma}[Convex Transform of a Submartingale]
Let $(Z_n)$ be a submartingale w.r.t.\ $(\mathcal{F}_n)$ and $\varphi:\mathbb{R}\to\mathbb{R}$ convex, nondecreasing, with $\mathbb{E}|\varphi(Z_n)|<\infty$ for all $n$. Then $(\varphi(Z_n))$ is a submartingale.
\end{lemma}

\begin{theorem}[Not-To-Gamble Theorem]
Let $(X_n)$ be a supermartingale and let $(H_n)$ be a predictable sequence (i.e.\ $H_n$ is $\mathcal{F}_{n-1}$-measurable) with $H_n\ge 0$ a.s. Define the stochastic integral
\[
(H\cdot X)_n := \sum_{i=1}^n H_i (X_i - X_{i-1}), \quad (H\cdot X)_0 := 0.
\]
Then $((H\cdot X)_n)$ is a supermartingale. Intuitively, no nonnegative predictable betting strategy can turn a losing game into a winning one in expectation.
\end{theorem}

\begin{theorem}[Optional Stopping for Sub/Supermartingales (Bounded Case)]
Let $(X_n)$ be adapted and $T$ a stopping time with $T\le N$ a.s.
\begin{itemize}
\item If $(X_n)$ is a submartingale, then $\mathbb{E}[X_N]\ge \mathbb{E}[X_T]\ge \mathbb{E}[X_0]$.
\item If $(X_n)$ is a supermartingale, then $\mathbb{E}[X_N]\le \mathbb{E}[X_T]\le \mathbb{E}[X_0]$.
\end{itemize}
\end{theorem}

\section{L8–L11: Martingale Convergence and Inequalities}

\begin{theorem}[Convergence of Nonnegative Supermartingales]
If $(X_n)$ is a supermartingale that is bounded below (e.g.\ $X_n\ge C$ a.s.), then there exists an integrable random variable $X_\infty$ such that $X_n\to X_\infty$ almost surely and in $L^1$ under mild conditions.
\end{theorem}

\begin{definition}[Doob's Maximal Inequality]
Let $(M_n)$ be a nonnegative submartingale and $a>0$. Then for any $n$,
\[
\mathbb{P}\Big(\max_{0\le k\le n} M_k \ge a\Big) \le \frac{\mathbb{E}[M_n]}{a}.
\]
This gives tail bounds for the running maximum of $(M_n)$.
\end{definition}

\begin{theorem}[Kolmogorov's Maximal Inequality]
Let $X_1,\dots,X_n$ be independent mean-zero random variables, and let $S_k = \sum_{i=1}^k X_i$. Then for any $a>0$,
\[
\mathbb{P}\Big(\max_{1\le k\le n} |S_k| \ge a\Big) \le \frac{\mathrm{Var}(S_n)}{a^2}.
\]
\end{theorem}

\begin{theorem}[Ville's Inequality]
Let $(M_n)$ be a nonnegative martingale with $\mathbb{E}[M_0]=1$. For any $a>0$,
\[
\mathbb{P}\Big(\sup_{n\ge 0} M_n \ge a\Big) \le \frac{1}{a}.
\]
This is a uniform-in-time version of Markov's inequality for martingales.
\end{theorem}

\begin{definition}[Subgaussian Random Variable]
A mean-zero random variable $X$ is \emph{subgaussian} with parameter $\sigma^2$ if
\[
\mathbb{E}[\exp(\lambda X)] \le \exp\left(\frac{\sigma^2\lambda^2}{2}\right) \quad \forall \lambda\in\mathbb{R}.
\]
Bounded variables are subgaussian with appropriate $\sigma^2$.
\end{definition}

\begin{theorem}[Hoeffding's Inequality]
Let $X_1,\dots,X_n$ be independent with $a_i\le X_i\le b_i$ a.s., and define $S_n = \sum X_i$, $\mu = \mathbb{E}S_n$. Then for all $t>0$,
\[
\mathbb{P}(S_n - \mu \ge t) \le \exp\left(-\frac{2t^2}{\sum_{i=1}^n (b_i-a_i)^2}\right),
\]
and similarly for the lower tail. This is a concentration inequality for bounded independent sums.
\end{theorem}

\begin{theorem}[Azuma–Hoeffding Inequality]
Let $(M_n)$ be a martingale with $M_0=0$, and assume bounded differences:
\[
|M_n - M_{n-1}|\le c_n \quad \text{a.s.\ for }n=1,\dots,N.
\]
Then for all $t>0$,
\[
\mathbb{P}(M_N \ge t) \le \exp\left(-\frac{2t^2}{\sum_{n=1}^N c_n^2}\right),
\]
and similarly for $\mathbb{P}(M_N \le -t)$. This is the martingale analogue of Hoeffding.
\end{theorem}

\begin{theorem}[McDiarmid's Inequality]
Let $X_1,\dots,X_n$ be independent random variables taking values in some set $\mathcal{X}$, and let $f:\mathcal{X}^n\to\mathbb{R}$ satisfy a bounded differences condition:
\[
\sup_{x_1,\dots,x_n,x_i'} |f(x_1,\dots,x_i,\dots,x_n) - f(x_1,\dots,x_i',\dots,x_n)| \le c_i.
\]
Let $Z = f(X_1,\dots,X_n)$ and $\mu = \mathbb{E}Z$. Then for all $t>0$,
\[
\mathbb{P}(Z-\mu \ge t) \le \exp\left(-\frac{2t^2}{\sum_{i=1}^n c_i^2}\right),
\]
and similarly for $\mathbb{P}(Z-\mu \le -t)$.
\end{theorem}

\section{L12–L15: Markov Chains, Stationarity, Convergence}

\begin{definition}[Discrete-Time Markov Chain]
A stochastic process $(X_n)_{n\ge 0}$ with state space $S$ is a \emph{Markov chain} with transition matrix $P$ if
\[
\mathbb{P}(X_{n+1}=y \mid X_0,\dots,X_n) = \mathbb{P}(X_{n+1}=y\mid X_n) = P(X_n,y)
\]
for all $n$ and states $y$. The $n$-step transition probabilities are $P^n(x,y) := \mathbb{P}_x(X_n=y)$.
\end{definition}

\begin{definition}[Stationary Distribution]
A probability vector $\pi$ on $S$ is \emph{stationary} for $P$ if
\[
\pi P = \pi.
\]
Equivalently, if $X_0\sim \pi$, then $X_n\sim \pi$ for all $n$.
\end{definition}

\begin{definition}[Reversibility and Detailed Balance]
A Markov chain with transition matrix $P$ and distribution $\pi$ is \emph{reversible} w.r.t.\ $\pi$ if
\[
\pi(x) P(x,y) = \pi(y) P(y,x) \quad \forall x,y\in S.
\]
If this holds, then $\pi$ is stationary. This is called the detailed balance condition.
\end{definition}

\begin{definition}[Irreducibility]
A Markov chain with state space $S$ is \emph{irreducible} if for all $x,y\in S$ there exists $n\ge 1$ such that $P^n(x,y)>0$. Intuitively, every state can reach every other state.
\end{definition}

\begin{definition}[Period and Aperiodicity]
The \emph{period} of a state $x$ is
\[
d(x) := \gcd\{n\ge 1 : P^n(x,x)>0\}.
\]
If $d(x)=1$, then $x$ is aperiodic. In an irreducible chain, all states have the same period $d$. The chain is \emph{aperiodic} if $d=1$.
\end{definition}

\begin{theorem}[Existence and Uniqueness of Stationary Distribution (Finite Case)]
Every finite irreducible Markov chain has at least one stationary distribution. In the finite irreducible case this stationary distribution is unique.
\end{theorem}

\begin{definition}[Total Variation Distance]
For two probability measures $\mu,\nu$ on a finite set,
\[
\|\mu - \nu\|_{\mathrm{TV}} = \frac12 \sum_{x} |\mu(x)-\nu(x)| = \sup_{A} |\mu(A)-\nu(A)|.
\]
\end{definition}

\begin{definition}[Mixing Time]
For a finite Markov chain with stationary distribution $\pi$, define
\[
d(t) := \sup_{x} \|P^t(x,\cdot) - \pi\|_{\mathrm{TV}}.
\]
Then the \emph{mixing time} is
\[
t_{\mathrm{mix}}(\varepsilon) := \min\{t: d(t)\le \varepsilon\}.
\]
\end{definition}

\begin{theorem}[Convergence Theorem]
For a finite, irreducible, aperiodic Markov chain with stationary distribution $\pi$,
\[
\lim_{n\to\infty} \|P^n(x,\cdot) - \pi\|_{\mathrm{TV}} = 0 \quad \text{for all }x.
\]
Moreover $\pi$ is unique.
\end{theorem}

\section{L16–L18: Reversible Chains, Spectrum, Spectral Gap}

\begin{definition}[Self-Adjointness and Spectrum for Reversible Chains]
For a reversible finite-state Markov chain with stationary $\pi$, consider the inner product
\[
\langle f,g\rangle_\pi := \sum_{x} f(x) g(x) \pi(x).
\]
Then $P$ is self-adjoint in $L^2(\pi)$:
\[
\langle Pf,g\rangle_\pi = \langle f,Pg\rangle_\pi.
\]
Hence all eigenvalues of $P$ are real and lie in $[-1,1]$, and there exists an orthonormal basis of eigenfunctions.
\end{definition}

\begin{definition}[Spectral Gap and Relaxation Time]
Let $1=\lambda_1 > \lambda_2\ge \dots \ge \lambda_n\ge -1$ be eigenvalues of $P$. Define
\[
\lambda^* := \max_{j\ge 2} |\lambda_j|, \qquad \gamma^* := 1-\lambda^*, \qquad t_{\mathrm{rel}} := \frac{1}{\gamma^*}.
\]
Here $\gamma^*$ is the \emph{absolute spectral gap} and $t_{\mathrm{rel}}$ is the \emph{relaxation time}.
\end{definition}

\begin{theorem}[Spectral Bound on Mixing Time]
For a finite, irreducible, reversible Markov chain with stationary distribution $\pi$,
\[
t_{\mathrm{mix}}(\varepsilon) \le t_{\mathrm{rel}} \log\left(\frac{1}{\varepsilon \,\pi_{\min}}\right),
\]
where $\pi_{\min} := \min_x \pi(x)$.
\end{theorem}

\begin{example}[Random Walk on an Undirected Graph]
Let $G=(V,E)$ be a connected undirected graph. Simple random walk from $x$ chooses a neighbor uniformly. Then
\[
\pi(x) = \frac{\deg(x)}{2|E|}
\]
is stationary and reversible. For regular graphs (e.g.\ cycles, hypercubes), $\pi$ is uniform.
\end{example}

\section{L19–L20: MCMC, Metropolis–Hastings, Gibbs}

\begin{definition}[Metropolis–Hastings Algorithm]
Let $\pi$ be a target distribution on finite state space $\Omega$, known up to a normalizing constant. Let $Q(x,y)$ be a proposal transition matrix. The MH chain has transition probabilities
\[
P(x,y) = Q(x,y)\alpha(x,y)\quad (y\neq x),
\]
where
\[
\alpha(x,y) = \min\left(1,\frac{\pi(y)Q(y,x)}{\pi(x)Q(x,y)}\right),
\]
and
\[
P(x,x) = 1 - \sum_{y\neq x} P(x,y).
\]
\end{definition}

\begin{theorem}[Detailed Balance for MH]
The Metropolis–Hastings chain satisfies
\[
\pi(x)P(x,y) = \pi(y)P(y,x)\quad \forall x,y.
\]
Hence $\pi$ is stationary, and the chain is reversible w.r.t.\ $\pi$.
\end{theorem}

\begin{definition}[Full Conditional and Gibbs Sampler]
Let $\pi$ be a joint distribution on $\Omega = \Omega_1\times\cdots\times\Omega_d$. The full conditional for coordinate $i$ is
\[
\pi(x_i \mid x_{-i}) := \mathbb{P}(X_i=x_i \mid X_j = x_j, j\neq i),
\]
viewed as a distribution on $\Omega_i$ with other coordinates fixed.

\emph{Random-scan Gibbs sampler}: at each step, choose coordinate $i$ (e.g.\ uniformly) and resample $X_i$ from $\pi(\cdot\mid X_{-i})$.
\end{definition}

\begin{lemma}[Proportional Trick for Conditionals]
If $\pi(x)\propto g(x)$ for some nonnegative function $g$, then
\[
\pi(x_i \mid x_{-i}) \propto g(x_i,x_{-i}) \quad \text{as a function of }x_i,
\]
with normalization over $x_i\in\Omega_i$.
\end{lemma}

\section{L21–L24: Coupling and Random Walk Recurrence}

\begin{definition}[Coupling of Probability Measures]
Let $\mu,\nu$ be probability distributions on the same state space $S$. A \emph{coupling} of $(\mu,\nu)$ is a pair of random variables $(X,Y)$ defined on a common probability space such that $X\sim \mu$ and $Y\sim \nu$.
\end{definition}

\begin{theorem}[Coupling Lemma]
For any two probability distributions $\mu,\nu$ on a finite set $S$,
\[
\|\mu-\nu\|_{\mathrm{TV}} = \inf\{ \mathbb{P}(X\neq Y) : (X,Y) \text{ is a coupling of }(\mu,\nu)\}.
\]
\end{theorem}

\begin{definition}[Markovian and Faithful Coupling]
Given a Markov chain with transition matrix $P$, a process $(X_t,Y_t)$ on $S\times S$ is a \emph{Markovian coupling} if each coordinate marginal evolves according to $P$:
\[
\mathbb{P}(X_{t+1}=x'\mid X_t=x,Y_t=y)=P(x,x'),\quad
\mathbb{P}(Y_{t+1}=y'\mid X_t=x,Y_t=y)=P(y,y').
\]
The coupling is \emph{faithful} if $X_t=Y_t$ implies $X_{t+1}=Y_{t+1}$ a.s. for all $t$.
\end{definition}

\begin{definition}[Coupling Time and Coupling Inequality]
For a coupling $(X_t,Y_t)$, the \emph{coupling time} is
\[
\tau := \inf\{t\ge 0 : X_t = Y_t\}.
\]
For any initial states $x,y$,
\[
\|P^t(x,\cdot) - P^t(y,\cdot)\|_{\mathrm{TV}} \le \mathbb{P}(X_t\neq Y_t) \le \mathbb{P}(\tau > t).
\]
In particular, if $Y_0\sim \pi$, then $\|P^t(x,\cdot)-\pi\|_{\mathrm{TV}} \le \mathbb{P}(\tau>t)$.
\end{definition}

\begin{definition}[Simple Random Walk on $\mathbb{Z}^d$]
Let $(X_n)$ be a Markov chain on $\mathbb{Z}^d$ where at each step the chain moves to a uniformly chosen neighbor (change one coordinate by $\pm 1$). This is SRW on $\mathbb{Z}^d$.
\end{definition}

\begin{theorem}[Recurrence and Transience of SRW]
Simple random walk on $\mathbb{Z}^d$ is:
\begin{itemize}
\item \emph{recurrent} for $d=1,2$ (returns to the origin infinitely often with probability 1),
\item \emph{transient} for $d\ge 3$ (returns only finitely many times with positive probability).
\end{itemize}
\end{theorem}


\section*{Canonical Worked Examples}

% =======================
% 1. Gambler's Ruin (Simple Symmetric RW)
% =======================

\subsection*{Example 1: Gambler's Ruin – Hitting Probability and Expected Time}

Consider simple symmetric random walk on $\{0,1,\dots,N\}$ with absorbing boundaries at $0$ and $N$ and $S_0=i$.

\paragraph{(a) Probability of hitting $N$ before $0$.}
Let
\[
p_i := \mathbb{P}_i(S_\tau = N),\quad \tau = \inf\{n\ge0: S_n\in\{0,N\}\}.
\]
By the Markov property and symmetry,
\[
p_i = \tfrac12 p_{i-1} + \tfrac12 p_{i+1},\quad 1\le i\le N-1,
\]
with boundary conditions $p_0=0$, $p_N=1$.
This is a linear second-order difference equation with general solution $p_i = a+bi$. Using $p_0=0$ gives $a=0$, and $p_N=1$ gives $b=1/N$. Thus
\[
\boxed{p_i = \frac{i}{N}.}
\]

\paragraph{(b) Expected time to absorption.}
Let
\[
h_i := \mathbb{E}_i[\tau].
\]
Again by conditioning on the first step,
\[
h_i = 1 + \tfrac12 h_{i-1} + \tfrac12 h_{i+1},\quad 1\le i\le N-1,
\]
with $h_0=h_N=0$. Rearranging:
\[
h_{i+1} - 2h_i + h_{i-1} = -2.
\]
A quadratic ansatz $h_i = \alpha i^2 + \beta i + \gamma$ yields
\[
h_{i+1}-2h_i+h_{i-1} = 2\alpha = -2 \Rightarrow \alpha=-1.
\]
Write $h_i = -i^2 + \beta i + \gamma$. Using $h_0=0$ gives $\gamma=0$, and $h_N=0$ gives $-N^2 + \beta N=0 \Rightarrow \beta = N$. Thus
\[
\boxed{h_i = i(N-i).}
\]

% =======================
% 2. Simple Random Walk as Martingale + Azuma
% =======================

\subsection*{Example 2: Simple Random Walk Martingale + Azuma}

Let $(X_i)$ be i.i.d.\ with $\mathbb{P}(X_i=1)=\mathbb{P}(X_i=-1)=\tfrac12$ and $S_n=\sum_{i=1}^n X_i$.

\paragraph{Martingale.}
With $\mathcal{F}_n=\sigma(X_1,\dots,X_n)$,
\[
\mathbb{E}[S_{n+1}\mid\mathcal{F}_n] = S_n + \mathbb{E}[X_{n+1}] = S_n,
\]
so $(S_n)$ is a martingale.

\paragraph{Azuma bound.}
Set $M_n = S_n$; then $|M_n - M_{n-1}| = |X_n|\le 1$. Azuma–Hoeffding gives, for all $t>0$,
\[
\mathbb{P}(S_n \ge t) \le \exp\!\Big(-\frac{2t^2}{\sum_{k=1}^n 1^2}\Big)
= \exp\!\Big(-\frac{2t^2}{n}\Big).
\]
Similarly for $\mathbb{P}(|S_n|\ge t)$ by a union bound.

% =======================
% 3. Doob Martingale + McDiarmid
% =======================

\subsection*{Example 3: Doob Martingale and McDiarmid}

Let $X_1,\dots,X_n$ be independent and $f:\mathcal{X}^n\to\mathbb{R}$ satisfy the bounded differences condition
\[
|f(x_1,\dots,x_i,\dots,x_n)-f(x_1,\dots,x_i',\dots,x_n)|\le c_i.
\]
Define $Z = f(X_1,\dots,X_n)$ and the Doob martingale
\[
M_k := \mathbb{E}[Z \mid X_1,\dots,X_k],\quad k=0,\dots,n.
\]
Then $(M_k)$ is a martingale and one can check $|M_k-M_{k-1}|\le c_k$ a.s. Applying Azuma,
\[
\mathbb{P}(Z - \mathbb{E}Z \ge t)
= \mathbb{P}(M_n - M_0 \ge t)
\le \exp\!\Big(-\frac{2t^2}{\sum_{k=1}^n c_k^2}\Big),
\]
which is McDiarmid's inequality.

% =======================
% 4. Stationary Distribution on an Undirected Graph
% =======================

\subsection*{Example 4: Random Walk on an Undirected Graph}

Let $G=(V,E)$ be a finite connected undirected graph. Simple random walk: from $x$, move to a uniformly chosen neighbor.

\paragraph{Stationary distribution.}
Define
\[
\pi(x) := \frac{\deg(x)}{2|E|}.
\]
Then for any adjacent $x,y$,
\[
\pi(x)P(x,y) = \frac{\deg(x)}{2|E|}\cdot \frac{1}{\deg(x)} = \frac{1}{2|E|}
= \frac{\deg(y)}{2|E|}\cdot \frac{1}{\deg(y)} = \pi(y)P(y,x),
\]
and both sides are $0$ if $(x,y)\notin E$. Thus detailed balance holds, so $\pi$ is stationary and the chain is reversible.

% =======================
% 5. Two-State Chain Coupling
% =======================

\subsection*{Example 5: Coupling for a Two-State Chain}

Consider $S=\{0,1\}$ and
\[
P = \begin{pmatrix}
3/4 & 1/4\\
1/2 & 1/2
\end{pmatrix}.
\]

\paragraph{Stationary distribution.}
Solve $\pi=\pi P$ with $\pi_0+\pi_1=1$:
\[
\pi_0 = \tfrac34\pi_0 + \tfrac12\pi_1 \Rightarrow \tfrac14\pi_0 = \tfrac12\pi_1 \Rightarrow \pi_1 = \tfrac12\pi_0.
\]
Normalization gives $\pi_0 + \tfrac12\pi_0 = 1 \Rightarrow \pi_0=2/3$, $\pi_1=1/3$.

\paragraph{Coupling construction.}
Define $r(0)=1/4,\ r(1)=1/2$. Let $U_{t+1}\sim\mathrm{Unif}(0,1)$ i.i.d., and set
\[
X_{t+1} = \mathbf{1}\{U_{t+1}\le r(X_t)\},\quad
Y_{t+1} = \mathbf{1}\{U_{t+1}\le r(Y_t)\}.
\]
Each marginal has transition $P$, and if $X_t=Y_t$ then $X_{t+1}=Y_{t+1}$ (faithful).

If $X_t\neq Y_t$ then one is $0$ and the other is $1$. Under the Bernoulli coupling, the mismatch probability is
\[
\mathbb{P}(X_{t+1}\neq Y_{t+1}\mid X_t\neq Y_t) = |r(1)-r(0)| = \tfrac14.
\]
Thus for $\tau=\inf\{t:X_t=Y_t\}$ and $X_0\neq Y_0$,
\[
\mathbb{P}(\tau>t)\le \Big(\tfrac14\Big)^t.
\]

\paragraph{Mixing bound.}
Start $X_0=x$ and $Y_0\sim\pi$. Then $Y_t\sim\pi$ and
\[
\|P^t(x,\cdot)-\pi\|_\mathrm{TV}\le \mathbb{P}(X_t\neq Y_t)\le \mathbb{P}(\tau>t)\le \Big(\tfrac14\Big)^t.
\]

% =======================
% 6. Lazy Random Walk on Cycle: Spectral Gap Order
% =======================

\subsection*{Example 6: Lazy RW on Cycle – Order of Mixing}

Lazy random walk on cycle of length $\ell$: at each step, stay w.p.\ $1/2$, move to each neighbor w.p.\ $1/4$.

\paragraph{Stationary distribution.}
Underlying graph is regular; transition matrix is symmetric; hence $\pi$ is uniform:
\[
\pi(x)=1/\ell.
\]

\paragraph{Spectral gap and mixing.}
Non-lazy RW on cycle has eigenvalues $\lambda_k = \cos(2\pi k/\ell)$, so $1-\lambda_2\asymp 1/\ell^2$. Laziness maps eigenvalues to $\lambda_k' = \tfrac12(1+\lambda_k)$, so the absolute spectral gap remains $\gamma^*\asymp 1/\ell^2$, and $t_\mathrm{rel}\asymp \ell^2$.

The spectral mixing bound gives
\[
t_\mathrm{mix}(\varepsilon)\ \lesssim\ t_\mathrm{rel}\,\log\!\Big(\frac{1}{\varepsilon\pi_{\min}}\Big)
\asymp \ell^2\log(\ell/\varepsilon),
\]
so for fixed $\varepsilon$, $t_\mathrm{mix}(\varepsilon)=O(\ell^2\log\ell)$.

% =======================
% 7. Metropolis--Hastings on 4 States
% =======================

\subsection*{Example 7: MH on $\{0,1,2,3\}$ with Path Proposal}

State space $\{0,1,2,3\}$, target $\pi(i)\propto i+1$ (weights $w(0)=1,w(1)=2,w(2)=3,w(3)=4$). Proposal: random walk on the path
\begin{itemize}
\item from $0$ propose $1$;
\item from $3$ propose $2$;
\item from $1,2$ propose each neighbor w.p.\ $1/2$.
\end{itemize}

\paragraph{Proposal kernel.}
Nonzero $q(i,j)$:
\[
q(0,1)=1,\ q(1,0)=q(1,2)=1/2,\ q(2,1)=q(2,3)=1/2,\ q(3,2)=1.
\]

\paragraph{Acceptance probabilities.}
For neighbors,
\[
\alpha(i,j)=\min\Big(1,\frac{w(j)q(j,i)}{w(i)q(i,j)}\Big).
\]
Compute e.g.\ from $1$:
\[
\alpha(1,0)=\min\Big(1,\frac{1\cdot 1}{2\cdot(1/2)}\Big)=1,\quad
\alpha(1,2)=\min\Big(1,\frac{3\cdot(1/2)}{2\cdot(1/2)}\Big)=1.
\]
Thus
\[
P(1,0)=q(1,0)\alpha(1,0)=\tfrac12,\quad
P(1,2)=q(1,2)\alpha(1,2)=\tfrac12,\quad P(1,1)=0.
\]
Similar calculations give all $P(i,j)$. Detailed balance with $\pi(i)\propto i+1$ holds by MH construction.

% =======================
% 8. Gibbs Sampler for 2-Spin Ising
% =======================

\subsection*{Example 8: Gibbs for 2-Node Ising Model}

Two spins $\sigma_1,\sigma_2\in\{-1,+1\}$ with joint distribution
\[
\pi(\sigma_1,\sigma_2)\propto \exp(\beta \sigma_1\sigma_2),\ \beta>0.
\]

\paragraph{Unnormalized weights.}
\[
(+1,+1): e^\beta;\quad (+1,-1): e^{-\beta};\quad (-1,+1): e^{-\beta};\quad (-1,-1): e^\beta.
\]

\paragraph{Full conditional for $\sigma_1$.}
Using proportional trick,
\[
\pi(\sigma_1\mid\sigma_2)\propto \exp(\beta \sigma_1\sigma_2).
\]
So if $\sigma_2=+1$,
\[
\mathbb{P}(\sigma_1=+1\mid\sigma_2=+1)=\frac{e^\beta}{e^\beta+e^{-\beta}}=\frac{1}{1+e^{-2\beta}},
\]
and if $\sigma_2=-1$,
\[
\mathbb{P}(\sigma_1=+1\mid\sigma_2=-1)=\frac{e^{-\beta}}{e^\beta+e^{-\beta}}=\frac{1}{1+e^{2\beta}}.
\]

\paragraph{Random-scan Gibbs step.}
At time $t$:
\begin{itemize}
\item With prob $1/2$, update $\sigma_1$ from $\pi(\cdot\mid \sigma_2)$.
\item With prob $1/2$, update $\sigma_2$ from $\pi(\cdot\mid \sigma_1)$ (symmetric formula).
\end{itemize}
This Gibbs chain is reversible w.r.t.\ $\pi$ and converges to $\pi$.

\section*{STAT 654 Cheatsheet (Key Facts)}

% --- Martingales, Stopping Times, Convergence ---

\textbf{Martingale:}
$(M_n,\mathcal{F}_n)$ with $\mathbb{E}|M_n|<\infty$ and $\mathbb{E}[M_{n+1}\mid\mathcal{F}_n]=M_n$ a.s.; sub/supermartingale if $\ge / \le$.  

\textbf{Stopping Time:}
$T$ s.t.\ $\{T\le n\}\in\mathcal{F}_n$ for all $n$; decision to stop uses only past info.  

\textbf{Stopped Process:}
$M_{n\wedge T}$ is again a martingale if $M_n$ is.  

\textbf{Optional Stopping (simple):}
If $(M_n)$ martingale and $T$ bounded a.s., then $\mathbb{E}[M_T]=\mathbb{E}[M_0]$.  

\textbf{Wald's Identity:}
i.i.d.\ $X_i$ with mean $\mu$, suitable stopping time $T$, then $\mathbb{E}\!\left[\sum_{i=1}^T X_i\right]=\mu\,\mathbb{E}[T]$.  

\textbf{Convex Transform:}
If $(M_n)$ martingale and $\varphi$ convex nondecreasing, then $(\varphi(M_n))$ is submartingale (when integrable).  

\textbf{Not-to-Gamble:}
If $(X_n)$ supermartingale and $H_n\ge0$ predictable, then $(\sum_{i=1}^n H_i(X_i-X_{i-1}))$ is supermartingale.  

\textbf{Supermartingale Convergence:}
Lower-bounded supermartingale converges a.s.\ (and in $L^1$ under mild conditions).  

% --- Big 3 Maximal / Tail Inequalities ---

\textbf{Doob Maximal Inequality:}
Nonnegative submartingale $(M_k)$, $a>0$,
\[
\mathbb{P}\!\left(\max_{0\le k\le n}M_k \ge a\right)\le \frac{\mathbb{E}[M_n]}{a}.
\]

\textbf{Kolmogorov Inequality:}
Independent mean-zero $X_i$, $S_k=\sum_{i=1}^k X_i$,
\[
\mathbb{P}\!\left(\max_{1\le k\le n}|S_k|\ge a\right)\le \frac{\mathrm{Var}(S_n)}{a^2}.
\]

\textbf{Ville Inequality:}
Nonnegative martingale $(M_n)$ with $\mathbb{E}[M_0]=1$,
\[
\mathbb{P}\!\left(\sup_{n\ge0}M_n\ge a\right)\le \frac{1}{a}.
\]

% --- Concentration: Hoeffding, Azuma, McDiarmid ---

\textbf{Hoeffding (bounded i.i.d.):}
$a_i\le X_i\le b_i$, $S_n=\sum X_i$, $\mu=\mathbb{E}S_n$,
\[
\mathbb{P}(S_n-\mu\ge t)\le \exp\!\Big(-\frac{2t^2}{\sum (b_i-a_i)^2}\Big).
\]

\textbf{Azuma–Hoeffding (martingale):}
$M_0=0$, $|M_n-M_{n-1}|\le c_n$,
\[
\mathbb{P}(M_N\ge t)\le \exp\!\Big(-\frac{2t^2}{\sum_{n=1}^N c_n^2}\Big).
\]

\textbf{McDiarmid (bounded differences):}
Independent $X_i$, $f$ with change $\le c_i$ per coordinate,
\[
\mathbb{P}(f(X)-\mathbb{E}f(X)\ge t)\le \exp\!\Big(-\frac{2t^2}{\sum c_i^2}\Big).
\]

% --- Gambler's Ruin / Random Walk ---

\textbf{Gambler's Ruin (simple RW on $\{0,\dots,N\}$):}
Start at $i$, $\tau=\inf\{n:S_n\in\{0,N\}\}$,
\[
\mathbb{P}_i(S_\tau=N)=\frac{i}{N},\quad \mathbb{E}_i[\tau]=i(N-i).
\]

\textbf{SRW Recurrence:}
SRW on $\mathbb{Z}^d$ recurrent for $d=1,2$, transient for $d\ge3$.  

% --- Markov Chains: Basic Definitions ---

\textbf{Markov Chain:}
$\mathbb{P}(X_{n+1}=y\mid X_0,\dots,X_n)=P(X_n,y)$; $P^n(x,y)=\mathbb{P}_x(X_n=y)$.  

\textbf{Stationary Distribution:}
$\pi$ with $\pi P=\pi$; if $X_0\sim\pi$ then $X_n\sim\pi$ for all $n$.  

\textbf{Reversible Chain:}
$\pi(x)P(x,y)=\pi(y)P(y,x)$ for all $x,y$ (detailed balance), implies $\pi$ stationary.  

\textbf{Irreducible:}
For all $x,y$, $\exists n: P^n(x,y)>0$.  

\textbf{Period / Aperiodic:}
$d(x)=\gcd\{n\ge1:P^n(x,x)>0\}$; chain aperiodic if $d(x)=1$ for all $x$ in class.  

\textbf{Convergence (finite, irr., aperiodic):}
Unique $\pi$, and $\|P^n(x,\cdot)-\pi\|_\mathrm{TV}\to0$ for all $x$.  

% --- Total Variation, Mixing Time, Coupling ---

\textbf{Total Variation:}
$\|\mu-\nu\|_\mathrm{TV}=\frac12\sum_x|\mu(x)-\nu(x)|=\sup_A|\mu(A)-\nu(A)|$.  

\textbf{Mixing Time:}
$d(t)=\sup_x\|P^t(x,\cdot)-\pi\|_\mathrm{TV}$,
\[
t_\mathrm{mix}(\varepsilon)=\min\{t:d(t)\le\varepsilon\}.
\]

\textbf{Coupling:}
$(X,Y)$ with $X\sim\mu$, $Y\sim\nu$; $\|\mu-\nu\|_\mathrm{TV}=\inf \mathbb{P}(X\neq Y)$.  

\textbf{Markovian Coupling:}
$(X_t,Y_t)$ on $S^2$ where each marginal evolves by $P$; faithful if $X_t=Y_t\Rightarrow X_{t+1}=Y_{t+1}$.  

\textbf{Coupling Time / Inequality:}
$\tau=\inf\{t:X_t=Y_t\}$,
\[
\|P^t(x,\cdot)-P^t(y,\cdot)\|_\mathrm{TV}\le\mathbb{P}(\tau>t).
\]

% --- Random Walk on Graphs ---

\textbf{Simple RW on Graph:}
From $x$, move to uniform neighbor; stationary $\pi(x)=\deg(x)/(2|E|)$ on connected undirected graphs.  

\textbf{Lazy RW:}
Stay w.p.\ $1/2$, otherwise move as original; kills period-2 issues on bipartite graphs.  

% --- Spectrum, Spectral Gap, Mixing Bounds ---

\textbf{Reversible Spectrum:}
$P$ self-adjoint in $L^2(\pi)$; eigenvalues real in $[-1,1]$.  

\textbf{Spectral Gap:}
$\lambda^*=\max_{j\ge2}|\lambda_j|$, $\gamma^*=1-\lambda^*$, $t_\mathrm{rel}=1/\gamma^*$.  

\textbf{Spectral Mixing Bound:}
\[
t_\mathrm{mix}(\varepsilon)\ \le\ t_\mathrm{rel}\,\log\!\left(\frac{1}{\varepsilon\,\pi_\mathrm{min}}\right),\quad
\pi_\mathrm{min}=\min_x\pi(x).
\]

% --- MCMC: Metropolis–Hastings and Gibbs ---

\textbf{Metropolis--Hastings:}
Proposal $Q(x,y)$, target $\pi$, accept prob
\[
\alpha(x,y)=\min\!\left(1,\frac{\pi(y)Q(y,x)}{\pi(x)Q(x,y)}\right).
\]
Transition: $P(x,y)=Q(x,y)\alpha(x,y)$ for $y\neq x$, $P(x,x)=1-\sum_{y\neq x}P(x,y)$.  

\textbf{MH Detailed Balance:}
$\pi(x)P(x,y)=\pi(y)P(y,x)$, so MH chain is reversible w.r.t.\ $\pi$.  

\textbf{Gibbs Sampler:}
On $\Omega=\Omega_1\times\cdots\times\Omega_d$, update coordinate $i$ by sampling from $\pi(x_i\mid x_{-i})$ (full conditional), leaving others fixed.  

\textbf{Proportional Trick:}
If $\pi(x)\propto g(x)$ then
\[
\pi(x_i\mid x_{-i})\propto g(x_i,x_{-i})
\]
as a function of $x_i$ (normalizing over $x_i$ only).  



\section*{Proof Techniques and Strategies (STAT 654)}



\subsection*{2. Martingale Techniques}

\textbf{Checking Martingale Property:}
Given candidate $(M_n)$, compute $\mathbb{E}[M_{n+1}\mid\mathcal{F}_n]$ explicitly.
Use linearity, independence of increments, and that conditioning removes future randomness.
Core examples: $S_n$, $S_n-n\mu$, $S_n^2-n\sigma^2$, exponential martingales.

\textbf{Constructing Martingales from Random Walks:}
Choose $f$ so that $f(S_n)$ has zero conditional drift:
\begin{itemize}
\item Linear: $S_n$, $S_n-n\mu$.
\item Quadratic: $S_n^2-n\sigma^2$.
\item Exponential: $r^{S_n} = \frac{q}{p}^{S_n  }$ with $r$ tuned to the step distribution (biased RW).
\end{itemize}
Then apply OST for hitting probabilities or expected times.

\textbf{Designing Martingales for Hitting Problems:}
Pick parameters to make boundary values convenient.
E.g.\ for biased RW on $\{0,\dots,N\}$, choose $M_n = \rho^{S_n}$ so that $\mathbb{E}[M_{n+1}\mid\mathcal{F}_n]=M_n$, and $M_\tau$ has easily computable values at $0$ and $N$.

\textbf{Careful Use of Optional Stopping:}
When using OST:
\begin{itemize}
\item Show $T$ is a stopping time.
\item Check conditions: $T$ bounded a.s., or integrability + bounded increments, etc.
\item Write $\mathbb{E}[M_T]=\mathbb{E}[M_0]$ and expand $M_T$ via terminal states (e.g.\ $0$ vs $N$).
\end{itemize}
This is the template for gambler’s ruin, biased RW, pattern waiting times.

\textbf{Applying Maximal Inequalities:}
Identify whether you have:
\begin{itemize}
\item Nonnegative submartingale $\Rightarrow$ Doob maximal inequality.
\item Independent mean-zero sums $\Rightarrow$ Kolmogorov inequality.
\item Nonnegative martingale with $\mathbb{E}[M_0]=1$ $\Rightarrow$ Ville inequality.
\end{itemize}
Verify assumptions, then plug into the inequality.

\textbf{Doob Martingale + Azuma/McDiarmid:}
Given $Z=f(X_1,\dots,X_n)$ with independent $X_i$:
\begin{itemize}
\item Define $M_k=\mathbb{E}[Z\mid X_1,\dots,X_k]$.
\item Show $|M_k-M_{k-1}|\le c_k$ (bounded differences).
\item Apply Azuma to $(M_k)$, which yields McDiarmid’s inequality for $f$.
\end{itemize}

\subsection*{3. Markov Chain and Random Walk Techniques}

\textbf{Solving $\pi P=\pi$:}
\begin{itemize}
\item Write linear equations for $\pi$ and use $\sum_x\pi(x)=1$.
\item Exploit symmetry (e.g.\ regular graphs, symmetric transitions).
\item Use degree rule: for simple RW on undirected graph, $\pi(x)\propto \deg(x)$.
\end{itemize}

\textbf{Checking Detailed Balance (Reversibility):}
For candidate $\pi$,
\[
\pi(x)P(x,y) \stackrel{?}{=} \pi(y)P(y,x).
\]
Check on edges / neighbor pairs; if holds and $P$ is supported on edges, then $\pi$ is stationary.

\textbf{Irreducibility Tests:}
Show from any $x$ you can reach any $y$ with positive probability in some number of steps:
\begin{itemize}
\item For RW on graph: show graph is connected.
\item For MCMC: show proposal and acceptance allow moving between all states.
\end{itemize}

\textbf{Aperiodicity Tests:}
\begin{itemize}
\item Show existence of self-loop $P(x,x)>0$ in irreducible chain $\Rightarrow$ all states aperiodic.
\item For cycles/paths: inspect possible return times to detect gcd of lengths.
\end{itemize}

\textbf{Graph-RW as Markov Chain:}
Convert graph to $P$:
\[
P(x,y)=\begin{cases}
1/\deg(x), & (x,y)\in E,\\
0, & \text{otherwise}.
\end{cases}
\]
Then apply stationary/reversibility facts: $\pi(x)=\deg(x)/(2|E|)$.

\textbf{Spectral Gap Reasoning:}
Use known behaviors:
\begin{itemize}
\item Cycle/path: gap $\asymp 1/n^2$ $\Rightarrow$ $t_{\mathrm{rel}}\asymp n^2$.
\item Complete graph: gap bounded away from 0 $\Rightarrow$ $t_{\mathrm{rel}}=O(1)$.
\item Hybercube: gap $\asymp 1/d$.
\end{itemize}
Then use
\[
t_{\mathrm{mix}}(\varepsilon)\lesssim t_{\mathrm{rel}}\log\frac{1}{\varepsilon\pi_{\min}}.
\]

\subsection*{4. Coupling Techniques}

\textbf{Shared-Randomness Markovian Coupling:}
Represent transitions from each state via intervals of $[0,1]$:
\begin{itemize}
\item For each state $i$, partition $[0,1]$ into subintervals of lengths $P(i,j)$.
\item At time $t$, draw a single $U_{t+1}\sim\mathrm{Unif}(0,1)$ and update both $X_{t+1}$ and $Y_{t+1}$ using their respective partitions for $X_t$ and $Y_t$.
\end{itemize}
This ensures each marginal evolves with kernel $P$.

\textbf{Faithful Coupling Design:}
Ensure that if $X_t=Y_t$ then the transition rule from that state is the same for both, using the same randomness, so $X_{t+1}=Y_{t+1}$ a.s.\ and once they meet they stay together.

\textbf{Geometric Tail Bounds for Coupling Time:}
If for all $i\neq j$,
\[
\mathbb{P}(X_{t+1}\neq Y_{t+1}\mid X_t=i,Y_t=j)\le c<1,
\]
then for $\tau=\inf\{t:X_t=Y_t\}$,
\[
\mathbb{P}(\tau>t)\le c^t.
\]
Technique: compute or bound one-step “stay different” probability and iterate.

\textbf{Distance-Based Couplings:}
Define a distance $d(X_t,Y_t)$ (graph distance, Hamming distance) and construct a coupling so $d$ evolves as a simpler process:
\begin{itemize}
\item Show $d$ has negative drift towards $0$ or behaves like a biased RW hitting $0$.
\item Use hitting-time estimates on $d$ to bound $\mathbb{P}(\tau>t)$.
\end{itemize}

\textbf{Coupling with Stationary Initial State:}
To bound $\|P^t(x,\cdot)-\pi\|_{\mathrm{TV}}$:
\begin{itemize}
\item Start $X_0=x$, $Y_0\sim\pi$.
\item Run coupling so $Y_t\sim\pi$ for all $t$.
\item Use $\|P^t(x,\cdot)-\pi\|_{\mathrm{TV}}\le \mathbb{P}(X_t\neq Y_t)\le \mathbb{P}(\tau>t)$.
\end{itemize}

\subsection*{5. MCMC (MH and Gibbs) Techniques}

\textbf{Designing MH Proposals:}
Choose $Q(x,y)$ so moves are:
\begin{itemize}
\item Local (random-walk proposals on graph or path).
\item Independence (always propose from a fixed distribution).
\item Symmetric ($Q(x,y)=Q(y,x)$) to simplify acceptance ratios.
\end{itemize}

\textbf{Computing MH Acceptances:}
Work with unnormalized target $w(x)\propto\pi(x)$:
\[
\alpha(x,y)=\min\Big(1,\frac{w(y)Q(y,x)}{w(x)Q(x,y)}\Big).
\]
Look for cancellations (especially if $Q$ symmetric) and compute off-diagonal $P(x,y)=Q(x,y)\alpha(x,y)$, then diagonals by row sums.

\textbf{MH Detailed Balance Proof:}
Show for $x\neq y$,
\[
\pi(x)Q(x,y)\alpha(x,y)=\pi(y)Q(y,x)\alpha(y,x),
\]
and note $P(x,x)$ defined by normalization gives detailed balance on diagonals as well.

\textbf{Gibbs Full-Conditional Derivation (Proportional Trick):}
If $\pi(x)\propto g(x)$, then for coordinate $i$:
\[
\pi(x_i\mid x_{-i})\propto g(x_i,x_{-i}) \quad \text{(treat $x_{-i}$ as fixed constant)}.
\]
Normalize over $x_i\in\Omega_i$ only.

\textbf{Gibbs Invariance/Stationarity:}
One-coordinate Gibbs update using full conditional $\pi(x_i\mid x_{-i})$ leaves $\pi$ invariant; compositions of such updates (systematic or random scan) also leave $\pi$ invariant, giving a reversible chain with stationary $\pi$.

\subsection*{6. Recurrence/Transience \& Distance Processes}

\textbf{Reduction to Known SRW Results:}
If a process can be expressed as a function of a random walk:
\begin{itemize}
\item Map to SRW or biased RW on $\mathbb{Z}$ or $\mathbb{Z}_{\ge0}$.
\item Use known formulas (gambler’s ruin, biased hitting probabilities) to infer recurrence, transience, or escape probabilities.
\end{itemize}

\textbf{Distance-from-Root Processes:}
On trees or graphs, define $L_n=\mathrm{dist}(X_n,\text{root})$:
\begin{itemize}
\item Show $L_n$ is a Markov chain with probabilities “up 2/3, down 1/3” etc.
\item Apply gambler’s ruin or biased RW analysis to $L_n$.
\end{itemize}

\subsection*{7. General Inequality and Bounding Tactics}

\textbf{Union Bound:}
For events $A_i$,
\[
\mathbb{P}\Big(\bigcup_i A_i\Big)\le \sum_i \mathbb{P}(A_i).
\]
Used with tail bounds to control maxima over indices.

\textbf{Jensen's Inequality:}
For convex $\varphi$ and integrable $X$,
\[
\varphi(\mathbb{E}[X\mid\mathcal{G}])\le \mathbb{E}[\varphi(X)\mid\mathcal{G}].
\]
Used in submartingale constructions and concentration proofs.

\textbf{Norm/Distance Comparisons:}
\begin{itemize}
\item Use coupling lemma: TV distance = inf disagreement probability over couplings.
\item Use Cauchy–Schwarz in $L^2(\pi)$ and spectral decomposition for eigenvalue-based bounds.
\end{itemize}

\end{document}